% ============================================================
% Model: 一维交替力常数链
% ============================================================

\section{一维交替力常数链}
\label{sec:1d-alternating-spring}

\begin{tipbox}[关键词标签]
交替力常数 · 相同原子 · 复式晶格 · 能带折叠 · 色散关系 · 态密度
\end{tipbox}

% --------------------------------------------------------
% 1. 模型定义框
% --------------------------------------------------------
\begin{modelbox}[模型：一维交替力常数链]
\textbf{物理情境}：一维链上所有原子质量相同（均为 $m$），但相邻原子间的弹簧劲度系数（力常数）交替为 $\alpha$ 和 $\beta$（设 $\alpha<\beta$）。这是一种``化学键交替''的一维模型，如聚乙炔的简化。

\textbf{运动方程}：设第 $2n$ 与 $2n+1$ 原子间力常数为 $\alpha$，第 $2n+1$ 与 $2n+2$ 原子间力常数为 $\beta$：
\begin{align*}
m\ddot{u}_{2n} &= \alpha(u_{2n+1}-u_{2n})+\beta(u_{2n-1}-u_{2n}),\\
m\ddot{u}_{2n+1} &= \beta(u_{2n+2}-u_{2n+1})+\alpha(u_{2n}-u_{2n+1}).
\end{align*}

\textbf{关键参数}：
\begin{itemize}[nosep]
    \item $m$：原子质量（全同）
    \item $\alpha, \beta$：交替力常数，$\alpha<\beta$
    \item $a$：相邻原子间距，晶格常数（原胞长度）$=2a$
    \item $\omega_0^2=(\alpha+\beta)/m$：特征频率平方
\end{itemize}

\textbf{边界条件 / 约束}：
\begin{itemize}[nosep]
    \item 仅近邻相互作用
    \item 周期性边界条件
    \item 第一布里渊区：$q\in[-\pi/2a, \pi/2a]$
\end{itemize}
\end{modelbox}

% --------------------------------------------------------
% 2. 关键公式框
% --------------------------------------------------------
\begin{formulabox}[核心公式]
\begin{enumerate}[label=\textbf{(\arabic*)}]
    \item \textbf{色散关系}：
    \[
    \omega_{\pm}^2=\frac{\alpha+\beta}{m}\pm\frac{1}{m}\sqrt{\alpha^2+\beta^2+2\alpha\beta\cos 2qa}
    \]
    \texteq{等价形式：$\omega_{\pm}^2=\frac{\alpha+\beta}{m}\Bigl[1\pm\sqrt{1-\frac{4\alpha\beta}{(\alpha+\beta)^2}\sin^2 qa}\Bigr]$}

    \item \textbf{布里渊区边界 ($q=\pm\pi/2a$)}：
    \[
    \omega_{-}=\sqrt{\frac{2\alpha}{m}},\qquad \omega_{+}=\sqrt{\frac{2\beta}{m}}
    \]
    \texteq{能隙：$\Delta\omega=\sqrt{\frac{2}{m}}(\sqrt{\beta}-\sqrt{\alpha})$}

    \item \textbf{长波极限 ($q\to0$)}：
    \[
    \omega_{-}\approx\sqrt{\frac{2\alpha\beta}{(\alpha+\beta)m}}\,|q|\cdot a,\qquad
    \omega_{+}\approx\sqrt{\frac{2(\alpha+\beta)}{m}}
    \]
    \texteq{声学支线性，光学支趋于常数}

    \item \textbf{$\alpha=\beta$ 退化}：
    \[
    \omega=2\sqrt{\frac{\beta_0}{m}}\Bigl|\sin\frac{qa}{2}\Bigr|,\qquad q\in\Bigl[-\frac{\pi}{a},\frac{\pi}{a}\Bigr]
    \]
    \texteq{退化为单原子链，原胞缩小为 $a$，布里渊区扩大一倍}

    \item \textbf{单原子链态密度}（$\alpha=\beta$ 时）：
    \[
    g(\omega)=\frac{2N}{\pi\sqrt{\omega_m^2-\omega^2}},\qquad \omega_m=2\sqrt{\frac{\beta_0}{m}}
    \]
\end{enumerate}
\end{formulabox}

% --------------------------------------------------------
% 3. 训练点框
% --------------------------------------------------------
\begin{drillbox}[训练点与考法变体]
\trainingpoint{1} \textbf{直接求解色散}：给定 $\alpha, \beta, m, a$，写运动方程、设行波解、推导久期方程并求解 $\omega_{\pm}(q)$。

\trainingpoint{2} \textbf{极限分析}：
\begin{itemize}[nosep]
    \item $\alpha\ll\beta$：声学支近似为``弱弹簧链''，光学支近似为``强弹簧振子''
    \item $\alpha=\beta$：严格证明退化为单原子链（能带折叠图像）
    \item $q\to0$：证明声学支声速 $v_s=a\sqrt{\frac{\alpha\beta}{(\alpha+\beta)m}}$
\end{itemize}

\trainingpoint{3} \textbf{能带折叠图像}：$\alpha=\beta$ 时，原双原子链的第一布里渊区 $[-\pi/2a,\pi/2a]$ 中两支色散，折叠进新的布里渊区 $[-\pi/a,\pi/a]$ 后连接成单支。

\trainingpoint{4} \textbf{与双原子链对比}：比较``质量不同+力常数相同''和``质量相同+力常数不同''两种模型的色散曲线异同。

\trainingpoint{5} \textbf{态密度作图}：给定色散关系，推导 $g(\omega)$，标出 van Hove 奇点位置。
\end{drillbox}

% --------------------------------------------------------
% 4. 解题技巧框
% --------------------------------------------------------
\begin{tipbox}[快速识别与解题技巧]
\begin{itemize}
    \item \examnote{识别标志}：题干出现``力常数交替''''相同原子''''不同弹簧''''聚乙炔型''，立即定位本模型。
    \item \examnote{久期方程速写法}：设 $u_{2n}=Ae^{i(2nqa-\omega t)}$，$u_{2n+1}=Be^{i[(2n+1)qa-\omega t]}$，直接写出
    \[
    \begin{pmatrix}\alpha+\beta-m\omega^2 & -(\alpha e^{iqa}+\beta e^{-iqa})\\ -(\beta e^{iqa}+\alpha e^{-iqa}) & \alpha+\beta-m\omega^2\end{pmatrix}\begin{pmatrix}A\\ B\end{pmatrix}=0.
    \]
    \item \examnote{非对角元对称性}：注意非对角元互为复共轭，$H_{12}=H_{21}^*$，保证久期方程为实数。
    \item \examnote{布里渊区判断}：周期为 $2a$（两个原子+两种弹簧构成一个重复单元），布里渊区边界在 $\pm\pi/2a$。
    \item \examnote{能隙估算}：$\alpha$ 与 $\beta$ 差异越大，能隙越大；$\alpha=\beta$ 时能隙闭合。
\end{itemize}
\end{tipbox}

% --------------------------------------------------------
% 5. 易错点框
% --------------------------------------------------------
\begin{errorbox}{常见失误与陷阱}
\begin{itemize}
    \item \textbf{运动方程符号}：$m\ddot{u}_{2n}=\alpha(u_{2n+1}-u_{2n})+\beta(u_{2n-1}-u_{2n})$ 中，$u_{2n-1}$ 与 $u_{2n}$ 之间是 $\beta$ 弹簧（因为 $2n-1$ 到 $2n$ 跨越了一个``奇--偶''边界）。不要误写成两个都是 $\alpha$。
    \item \textbf{非对角元展开}：$\alpha e^{iqa}+\beta e^{-iqa}=(\alpha+\beta)\cos qa+i(\alpha-\beta)\sin qa$，其实部决定两支的``耦合强度''。
    \item \textbf{退化验证}：$\alpha=\beta$ 时，$\omega_{\pm}^2=\frac{2\alpha}{m}(1\pm\cos qa)$，即 $\omega_{-}=\sqrt{\frac{2\alpha}{m}}|\sin(qa/2)|$，$\omega_{+}=\sqrt{\frac{2\alpha}{m}}|\cos(qa/2)|$。必须将 $q$ 范围扩展到 $[-\pi/a,\pi/a]$ 才能看到单支。
    \item \textbf{态密度积分}：$\int_0^{\omega_m}g(\omega)\,d\omega=N$，若结果不是 $N$ 说明计算有误。
    \item \textbf{与标准双原子链混淆}：标准双原子链是``质量不同+力常数相同''，这里是``质量相同+力常数不同''。运动方程的结构不同，不要套用旧公式。
\end{itemize}
\end{errorbox}

% --------------------------------------------------------
% 6. 物理图像与关联模型
% --------------------------------------------------------
\begin{insightbox}[物理图像与模型关联]
\textbf{物理图像}：

想象一条由相同小球串成的项链，但项链的链子交替使用``橡皮筋''（$\alpha$）和``钢丝''（$\beta$）：
\begin{itemize}[nosep]
    \item \textbf{声学支}：整体晃动，橡皮筋和钢丝一起伸缩。长波时几乎不拉伸任何弹簧，频率趋于零。
    \item \textbf{光学支}：相邻小球反向振动，橡皮筋被剧烈拉伸/压缩而钢丝几乎不动（若 $\beta\gg\alpha$）。振动频率由较软的弹簧决定。
\end{itemize}

\textbf{与相似模型的对比}：
\begin{center}
\begin{tabular}{llll}
\toprule
\textbf{对比维度} & \textbf{单原子链} & \textbf{双原子链} & \textbf{交替力常数链} \\
\midrule
原子质量 & 全同 $m$ & $M\neq m$ & 全同 $m$ \\
力常数 & 全同 $\beta$ & 全同 $\beta$ & 交替 $\alpha\neq\beta$ \\
原胞长度 & $a$ & $2a$ & $2a$ \\
声学支 & 有 & 有 & 有 \\
光学支 & 无 & 有 & 有 \\
退化条件 & --- & $M=m$ 时退化为单原子链 & $\alpha=\beta$ 时退化为单原子链 \\
\bottomrule
\end{tabular}
\end{center}

\textbf{推广方向}：
\begin{itemize}[nosep]
    \item 聚乙炔（SSH 模型）：电子--晶格耦合导致键长交替（Peierls 不稳定性），打开能隙，形成拓扑孤子
    \item 更复杂的周期：力常数按 $\alpha,\beta,\gamma$ 周期排列，出现更多支
    \item 无序链：Anderson 局域化，所有本征态指数衰减
\end{itemize}
\end{insightbox}

% --------------------------------------------------------
% 7. 推导附录
% --------------------------------------------------------
\begin{derivationbox}[推导细节（自我检验用）]
\textbf{Step 1：设行波解}

设 $u_{2n}=Ae^{i(2nqa-\omega t)}$，$u_{2n+1}=Be^{i[(2n+1)qa-\omega t]}$。

\textbf{Step 2：代入运动方程}

\begin{align*}
-m\omega^2 A &= \alpha(Be^{iqa}-A)+\beta(Be^{-iqa}-A)=(\alpha+\beta)(B\cos qa-A)+i(\alpha-\beta)B\sin qa,\\
-m\omega^2 B &= \beta(Ae^{iqa}-B)+\alpha(Ae^{-iqa}-B)=(\alpha+\beta)(A\cos qa-B)+i(\alpha-\beta)A\sin qa.
\end{align*}

整理为矩阵：
\[
\begin{pmatrix}\alpha+\beta-m\omega^2 & -(\alpha e^{iqa}+\beta e^{-iqa})\\ -(\beta e^{iqa}+\alpha e^{-iqa}) & \alpha+\beta-m\omega^2\end{pmatrix}\begin{pmatrix}A\\ B\end{pmatrix}=0.
\]

\textbf{Step 3：久期方程}

$(\alpha+\beta-m\omega^2)^2=|\alpha e^{iqa}+\beta e^{-iqa}|^2=\alpha^2+\beta^2+2\alpha\beta\cos 2qa$。

故
\[
\omega_{\pm}^2=\frac{\alpha+\beta}{m}\pm\frac{1}{m}\sqrt{\alpha^2+\beta^2+2\alpha\beta\cos 2qa}.
\]

\textbf{Step 4：$\alpha=\beta$ 退化}

$\omega_{\pm}^2=\frac{2\alpha}{m}\bigl(1\pm\cos qa\bigr)$。

$\omega_{-}=\sqrt{\frac{2\alpha}{m}}|\sin(qa/2)|$，$\omega_{+}=\sqrt{\frac{2\alpha}{m}}|\cos(qa/2)|$。

将 $q$ 从 $[-\pi/2a,\pi/2a]$ 扩展到 $[-\pi/a,\pi/a]$，$\omega_{+}(q)=\omega_{-}(q+\pi/a)$，两支连接成单原子链色散。

\textbf{Step 5：态密度}

$\alpha=\beta$ 时退化为标准单原子链，$g(\omega)=\frac{2N}{\pi\sqrt{\omega_m^2-\omega^2}}$，$\omega_m=2\sqrt{\alpha/m}$。
\end{derivationbox}

\vspace{1em}
